\section{Historical Data: Coverage and Available Records}
\label{app:archive-scope}

The archived dataset contains 1,028,822 raw records. After excluding unsuccessful responses and removing duplicate texts, we obtain 899,025 model--benchmark--item--response records when
recorded input variants remain distinct, and 837,704 when variants are ignored
for text inventory only. Neither count estimates independent generation events.
Table~\ref{tab:archive-roles} classifies all 39 benchmark settings by reviewing the current benchmark implementations. Each classification
records the relevant source file, class, line, hash, and reasoning. These
implementations help identify task roles, but do not establish the exact
prompts used for every historical request.

\begin{table}[ht]
\centering\small
\begin{tabular}{lrr}
\toprule
Current task role & Settings & Records\\
\midrule
Evaluation of existing work & 7 & 354,308\\
Test or option answering & 10 & 306,307\\
Mixed educational assistance & 2 & 78,322\\
Tutor response & 9 & 68,799\\
Learner history or diagnosis & 2 & 42,997\\
Safety response & 2 & 13,974\\
Metacognitive problem solving & 3 & 11,188\\
Constrained question generation & 1 & 9,233\\
Assigned solution repair & 1 & 8,015\\
General instruction following & 1 & 3,782\\
Knowledge-structure reasoning & 1 & 2,100\\
\midrule
Total & 39 & 899,025\\
\bottomrule
\end{tabular}
\caption{Counts of distinct texts with input versions kept separate, not counts of
independent tutoring interactions or personality labels. Mixed educational assistance requires item-level
role separation; it is not automatically excluded from later research.}
\label{tab:archive-roles}
\end{table}

Evaluation and test/option outputs comprise 73.48\% of these records. For
example, the Pedagogy Benchmark requests pedagogical-knowledge multiple-choice
answers, while Automated Student Assessment Prize (ASAP) and Short Answer
Scoring (SAS) settings ask the candidate model to score student work.
MRBench and Building Educational Applications (BEA) judge settings likewise
evaluate existing tutor text; that text
is not the candidate evaluator's own tutoring. Even among tutor settings,
instructions differ: TutorBench use cases request different actions, MMTutorBench
requires image-conditioned guidance in three prescribed sections, and LongTutor
\citep{li2026longtutor}
explicitly requests scaffolding without the full answer.

A second pass verifies all 122 prediction files and corresponding summaries in
the nine tutor-response settings against the census hashes. Their 86,072 latest
successful file items reproduce 68,799 variant-preserving deduplicated records.
None of those file items contains the inspected fields for complete messages,
prompt text or input hashes, provider request identifiers, returned model names,
generation timestamps, or finish reasons. All contain usage metadata. Of the
122 summaries, 83 contain generation parameters and 114 contain a judge-prompt
hash; none contains the inspected generation-prompt hash fields. A judge-prompt
hash is not evidence of the tutor's generation prompt, and later scoring summaries
do not necessarily establish the original generation parameters.

These fields are absent from the inspected files; other logs or original
sources may still contain them. Reconstructed historical inputs must be
identified as reconstructions. We therefore use the archive to identify candidate patterns and investigate
alternative explanations. The controlled study records inputs, outputs,
problem sources, model versions, and timing for its more focused tests.

\subsection{What the Existing Scores Contribute}
\label{app:archive-reanalysis}

We also reanalyze previously scored responses from six models on the original
eight behavioral dimensions. These are exploratory measurements,
not labels assigned to the entire archive, and are not pooled with the new
binary event codes. Excluding 15 contexts without recoverable source groups
leaves 224 contexts from 210 source families across four task settings. Each
source stays in one fold; contexts share their source's weight within a task,
and tasks receive equal weight. Across 100 deterministic two-fold splits, we
compare a global model profile with a separate model profile for each known
task, predicting scores centered within each context. This tests new sources
within known task settings, not transfer to a new task or prediction of absolute
scores.

The gain from task-specific profiles depends on which instructional settings are
included (Table~\ref{tab:archive-profile-sensitivity}). Excluding the Socratic
setting, whose current adapter requires question-only output, reduces the
direct-help increment from 6.57\% to 0.34\%. This does not eliminate predictive
information in the global model profile: in the reduced dataset, its direct-help
mean squared error (MSE) is 0.66754, compared with 0.88413 for zero relative model differences.
Personalization retains a larger gain from task-specific profiles, but that legacy
score does not establish learner understanding or a validated personality factor.

\begin{table}[ht]
\centering\small
\begin{tabular}{lrr}
\toprule
Legacy dimension & Four tasks & Without Socratic\\
\midrule
Affective warmth & 11.27\% & 2.81\%\\
Autonomy support & 9.58\% & 4.49\%\\
Cognitive load & 2.05\% & 1.83\%\\
Diagnostic specificity & 3.93\% & 2.51\%\\
Elicitation & 10.73\% & 3.23\%\\
Epistemic caution & $-0.86$\% & $-1.11$\%\\
Help directness & 6.57\% & 0.34\%\\
Personalization & 12.15\% & 10.52\%\\
\bottomrule
\end{tabular}
\caption{Exploratory relative MSE reduction from task-specific model profiles
over a global model profile, averaged over split-specific ratios. The full
dataset has 224 contexts/210 source families; the reduced dataset has 144/133.
Negative values mean worse predictions. The 100 overlapping splits are not
independent experiments and their dispersion is not a confidence interval.}
\label{tab:archive-profile-sensitivity}
\end{table}

This exclusion was chosen after observing task results. It changes the target
task distribution and its weights, so the difference between columns is not a
causal effect of question-only formatting. We retain both analyses. Their role
is to motivate separating problem content, student cues, and output instructions
in the prospective experiment; neither archive size nor favorable historical scores replaces those tests.

\subsection{Data Use, Access, and Privacy}
\label{app:data-use}
The prospective problems and student messages are synthetic English research
stimuli, not a demographic sample. Archived tasks have heterogeneous languages
and roles; the million-record census is not a uniformly English tutoring
corpus. Existing human teacher-action labels and LongTutor diagnosis labels
are secondary reference data, not new human annotations collected here.

The MathDial repository \citep{macina2023mathdial} states Creative Commons
Attribution--ShareAlike (CC BY-SA) 4.0.\footnote{\url{https://github.com/eth-nlped/mathdial}}
For MathTutorBench \citep{macina2025mathtutorbench}, the publication states
Creative Commons Attribution (CC BY) 4.0, whereas the current repository README
states CC BY-SA 4.0.\footnote{\url{https://github.com/eth-lre/mathtutorbench}}
We retain this discrepancy rather than infer a uniform redistribution licence.
LongTutor's data card \citep{li2026longtutor} states CC BY 4.0 for data and MIT
for evaluation code. These statements identify source terms, not independent
verification of every underlying record's permissions. Archived generations
also retain the conditions of their upstream sources and application programming
interface (API) services.

Raw archived learner histories and the full evaluation corpus are not
redistributed with this manuscript. Earlier semantic payload preparation used
identifier-pattern screening, which cannot establish absence of all sensitive
information. The appendix instead provides aggregate results, measurement
rules, synthetic stimuli, and worked response examples from the prospective
experiment. The intended use is research on model behavior, not reconstruction
of learners or inference of their personal traits. No institutional ethics
approval or new participant consent is asserted for this secondary analysis.
